\runninghead{Each Text in Its Setting}
\section{Each Text in Its Setting}
\label{sec:each-element}

Each appointed text is set here in its own place: the chapter or psalm it
comes from, the words the Missal keeps and leaves out, its earlier use where
that is known, and what the Fathers and later commentators say of its words.
The three interpretations that follow take up the larger questions.

\subsection{Introit: the Church's antiphon and the psalm of the just}
\label{sec:element-introit}

The antiphon was not written for this feast. Blessed Ildefonso Schuster notes
that \latin{Gaudeamus} is an ancient Introit first assigned to the feast of St
Agatha on 5~February (\work{The Sacramentary} V, p.~209), and William Durandus,
in the thirteenth century, gives \latin{Gaudeamus omnes in Domino} to St
Agatha (\work{Rationale} VII.47). On this feast the antiphon names, in place of
one martyr, \latin{Sanctorum omnium}, and keeps the angels who rejoice at the
solemnity and praise the Son of God. The Gregorian Antiphonary printed by
Migne already gives the Introit \latin{Gaudeamus omnes in Domino} for All
Saints, with the verse \latin{Exsultate justi in Domino, rectos decet
collaudatio}. Lucien Fromage, continuing Dom Prosper Guéranger's
\work{L'Année liturgique}, comments on the antiphon's angels; his comment is
given with the first interpretation below.

The verse opens Psalm~32 (33), a hymn in which the just are bidden to rejoice
and to praise on the harp, and which goes on to the word and works of the Lord
and his counsel that stands for ever. St Augustine reads the first verse as a
warning against a joy that rests on oneself: \latin{exsultate justi, non in
vobis, non enim hoc tutum est; sed in Domino}, ``rejoice, you just, not in
yourselves, for that is not safe, but in the Lord''; and only those praise the
Lord \latin{qui se subdunt Domino}, ``who submit themselves to the Lord''
(\work{Enarr.\ in Ps.}\ 32, en.~1). In his second exposition he asks who the
upright are, and answers, \latin{Qui dirigunt cor secundum voluntatem Dei; et
si eos perturbat humana fragilitas, divina consolatur aequitas}: those who
direct the heart by the will of God, and whom, if human frailty troubles them,
God's justice consoles (en.~2, s.~1).

St Basil the Great tells the just to exult \latin{non cum res domi vobis
prospere succedunt \dots\ sed quod Dominum habetis}, not when things go well
at home, but because they have the Lord (\work{Hom.\ in Ps.}\ 32, n.~1). St
Thomas Aquinas draws a meaning from the prefix of \latin{collaudatio}: those
whose will is conformed to God's \latin{gaudent in prosperis et adversis; et
ideo dicit Collaudatio, quia de omnibus laudant, non in aliquibus tantum. Item
unanimiter}, ``rejoice in prosperity and adversity, and so he says
\latin{collaudatio}, because they praise in all things, not in some only;
and also with one mind'' (\work{In Ps.}\ 32). St Robert Bellarmine reads the
second half of the verse as a special invitation to the just, with whom God is
always pleased. An extract under the name of Cassiodorus in St Bruno of
Würzburg's commentary on the Psalms says the same of the prefix in another
way: \latin{Collaudatio vero multorum ore una laudatio est, propter unitatem
Ecclesiae}, ``praise together is one praise from the mouths of many, because of
the unity of the Church.''

\subsection{Collect: one celebration and many intercessors}
\label{sec:element-collect}

The Collect is older than the feast. The Leonine Sacramentary, in its prayers
for the feast of Ss Peter and Paul in June, and the Gelasian Sacramentary,
among its evening prayers after the same Masses, give it as a prayer for all
the apostles: \latin{qui nos omnium apostolorum merita sub una tribuisti
celebritate venerari}. Neither book has a Mass of All Saints. Where they ask
for the abundance of mercy \latin{celerem} or \latin{celeriter}, ``swift'' or
``swiftly,'' the feast's Collect reads \latin{desideratam}, ``desired,'' the
reading that Wilson's edition reports from a later hand of a St Gall
manuscript and from Gerbert's edition. Schuster says the Collect is Gelasian and belonged to a
collective feast of the Apostles (p.~210).

The Collect's last words exist in two forms. The apostles' prayer in both
sacramentaries, the Gilbertine and Sarum Missals, and the Roman Missal from
1570 to 1962 read \latin{multiplicatis intercessoribus}, ``the intercessors
being multiplied.'' The Gregorian Sacramentary in Hugues Ménard's codex, the
eleventh-century sacramentary of Monte Suavicino, and Durandus's
\work{Rationale} read \latin{multiplicatis intercessionibus}, ``the
intercessions being multiplied.'' The age of the books does not show which
reading came first, since the Leonine and Gelasian books give the apostles'
prayer and not its adaptation to all the saints. Ménard says only that his
codex's reading, though the Roman Missal reads \latin{intercessoribus},
\latin{non est contemnenda}, ``is not to be despised.''

Both forms make the petition rest on many asking. Durandus gives four causes
for which the feast was instituted: the dedication of the Pantheon; the
supplying of the feasts of saints who have none; the making good of
negligence in keeping them; and the easier granting of petitions, for which
he quotes this Collect (\work{Rationale} VII.34; his reason is given in the
first interpretation below). St Bernard, in a sermon for the feast, says that our infirmity is helped by the saints'
intercession, \latin{cum eorum intercessione juvetur infirmitas nostra}
(\work{In festo omnium Sanctorum}, sermo~2, n.~1).

\subsection{Commemoration Collect: the bonds of sins}
\label{sec:element-commemoration-collect}

The Sunday's Collect is said because the feast has displaced the Sunday, after
the feast's Collect and under its own conclusion. It asks that God absolve the
sins of his peoples, so that by his kindness they may be freed from the bonds
that their frailty has contracted. It was composed for the Twenty-third Sunday
after Pentecost, not for this feast, and no commentator on the feast speaks of
it. In the order in which the priest says them, the feast's prayer asks for
mercy through the many saints and the Sunday's asks release from the sins of
those who pray.

\subsection{Lesson: the sealing and the multitude}
\label{sec:element-lesson}

The seventh chapter of the Apocalypse stands between the opening of the sixth
seal and the seventh. Four angels stand at the four corners of the earth,
holding the four winds (7:1). The angel from the rising of the sun forbids
them to hurt the earth or the sea or the trees until the servants of God are
sealed on their foreheads (7:2--3). John hears the number, a hundred and
forty-four thousand, twelve thousand from each tribe, and then sees a
multitude ``which no man could number,'' from all nations, in white robes and
with palms. The angels, the ancients and the four living creatures stand round
the throne, and the angels fall on their faces and adore God with a sevenfold
praise. The Lesson stops before the elder's question in 7:13 and his answer,
that these ``are come out of great tribulation'' (7:14).

Schuster takes the Lesson as a sign of the feast's origin, since the same
Lesson served for the dedication of the Pantheon on 13~May (p.~210; he prints
chapter~viii for vii). Durandus gives the same Lesson, \latin{Ecce ego
Ioannes}, in his account of the feast, and Sicard of Cremona says why it suits
the day: \latin{Huic festo bene congruit epistola de numero signatorum, et evangelium de numero beatitudinum}, ``to this feast
there fits well the epistle on the number of the sealed and the gospel on the
number of the beatitudes'' (\work{Mitrale} IX.49). The Fathers' answers to
who the sealed and the multitude are belong to the first interpretation
below.

\subsection{Gradual: fear and want}
\label{sec:element-gradual}

Psalm~33 (34) is an alphabetical psalm. Its title sets it at the time when
David ``changed his countenance before Achimelech, who dismissed him, and he
went his way'' (33:1; compare 1~Kings 21:10--15). In the verses before the
Gradual's, the angel of the Lord encamps round those who fear him and
delivers them, and the psalm says, ``O taste, and see that the Lord is sweet''
(33:8--9). The Gradual takes up at the next verse, ``Fear the Lord, all ye his
saints,'' and leaves out the first half of verse~11, ``The rich have wanted,
and have suffered hunger.'' The Gregorian Antiphonary already gives the
Gradual for All Saints.

The psalm's own word for the faithful, \latin{Sancti eius}, ``his saints,'' is
the word of the feast. St Bernard, preaching on the feast, speaks of the verse
as it was sung in the night office, and his sermon is given in the third
interpretation below. The promise attached to the verse, \latin{nihil deest
timentibus eum}, meets the objection St Augustine puts in the mouth of a
man who seeks the Lord daily and has nothing good, while the wicked rich die
in their riches; Augustine, St Basil, St Thomas and St Robert Bellarmine
answer it in the third interpretation below.

\subsection{Alleluia: the invitation}
\label{sec:element-alleluia}

The verse is the Lord's invitation in Mt 11:28. Before it he has upbraided
Corozain, Bethsaida and Capharnaum (11:20--24), and thanked the Father for
hiding ``these things from the wise and prudent'' and revealing them ``to
little ones'' (11:25--27); the place where he speaks is not named. After it he
says, ``Take up my yoke upon you, and learn of me, because I am meek, and
humble of heart \dots\ For my yoke is sweet and my burden light''
(11:29--30). Schuster calls the verse a preparation for the Gospel (p.~210).

The Fathers differ on what burden the invitation lifts. St John Chrysostom
says it calls ``not this or that person, but all that are in anxiety, in
sorrows, in sins,'' and that ``sin too hath labor, and a burden''
(\work{Hom.\ in Matt.}\ 38.3). St Augustine, preaching on the verse, says that
all labour ``except that we are mortal men, frail creatures and infirm,
bearing about vessels of clay,'' and asks, ``What then does He mean by, `Come
unto Me, all ye that labour,' but that ye may labour no more?'' (\work{Sermon}
XIX [LXIX], n.~1). St Hilary of Poitiers reads the labour as the difficulties
of the Law and the burden as the world's sins
(\latin{Legis deinde difficultatibus laborantes, et peccatis saeculi
oneratos}), and says the Lord promises to take them away if they take his
yoke, that is, his commandments, and come to him \latin{sacramento crucis},
by the sacrament of the cross (\work{In Matt.}\ XI, n.~15). St Jerome gives
two answers and does not choose: the burden is sin, as Zacharias and the
psalmist testify (Zach 5; Ps 37 (38)), \latin{vel certe}, ``or else,'' the
heavy yoke of the Law, from which the Lord calls to the grace of the Gospel
(\work{In Matt.}\ II); Rabanus Maurus takes the comment over nearly word for
word in his own commentary. St Gregory the Great applies the verse to those of
Judaea who ``were weighed down with the heavy yoke of the Law'' (\work{Moralia}
I, n.~24), and St Thomas reports Hilary's sentence to show that the New Law is
a lighter burden than the Old (\work{ST} I-II q.~107 a.~4).

\subsection{Gospel: the Beatitudes}
\label{sec:element-gospel}

The Gospel is the opening of the Sermon on the Mount. Jesus has gone ``about
all Galilee,'' and crowds have followed him from Galilee, Decapolis,
Jerusalem, Judea and beyond the Jordan (4:23--25). Seeing the multitudes, he
goes up into a mountain that Matthew does not name, sits, and teaches his
disciples (5:1--2). The eight beatitudes follow (5:3--10), and then he applies
the last of them to the disciples themselves: they are blessed when they are
reviled and persecuted for his sake, and their reward is very great in heaven
(5:11--12a). The Missal stops before ``For so they persecuted the prophets
that were before you.'' The first and the eighth beatitude end with the same
promise, ``for theirs is the kingdom of heaven,'' in the present tense; the six
between promise in the future.

Schuster says that on this feast the Gospel could only be that of the
Beatitudes, since all are included in it, and that ``those who have least
according to the world receive a larger share of the heavenly gift''
(pp.~210--211). Durandus and Sicard give the same Gospel, and Sicard sets its
number beside the Lesson's. St Bernard's first sermon for the feast is an
exposition of these verses. The second interpretation below takes up Bernard
and the Fathers on the order of the beatitudes and the time of their
rewards.

\subsection{Offertory: in the hand of God}
\label{sec:element-offertory}

The Book of Wisdom gives first the speech of the wicked, who reckon life short
and death the end, and resolve, ``Let us oppress the poor just man \dots\ Let
us condemn him to a most shameful death'' (Wis 2:10, 20). Chapter~3 answers
them: the souls of the just are in the hand of God; to the unwise they seemed
to die, but they are in peace; ``their hope is full of immortality'' (3:4).
The antiphon keeps the first clause, the appearance of death and the peace,
and leaves out the unwise men's judgement that the departure of the just was
misery and destruction.

Schuster reads the antiphon of the martyrs, and St Alphonsus Liguori of every
just soul; St Augustine and St Bonaventure gloss its first clause; their
readings are given in the third interpretation below. Fromage connects the
antiphon to the Gospel just read, as the second interpretation sets out.
Other Fathers use the verse without expounding it. St John
Chrysostom cites ``the souls of the righteous are in the hand of God'' to deny
that the souls of the dead wander the earth (\work{Hom.\ in Matt.}\ 28, n.~3),
and St Gregory the Great applies the following clause to the martyr whose
body is bereft of life ``amidst the hands of the persecutors'' (\work{Moralia}
XII, n.~6).

\subsection{Secret: the honour of all the just}
\label{sec:element-secret}

The Secret offers ``the gifts of our devotion'' and asks two things of them:
that they be pleasing to God \latin{pro cunctorum \dots\ honore iustorum}, for
the honour of all the just, and that by his mercy they be for our salvation.
The saints are named here by the word the Offertory has just sung,
\latin{iustorum}. The prayer stands in the Monte Suavicino sacramentary with
the same word; Ménard's Gregorian Sacramentary reads \latin{honore sanctorum},
``the honour of the saints.'' Schuster translates the prayer and adds nothing to
it (p.~211). It asks that the offering honour the saints and benefit those who
make it, and it asks nothing for the saints themselves.

\subsection{Commemoration Secret: the sacrifice of praise}
\label{sec:element-commemoration-secret}

The Sunday's Secret offers \latin{sacrificium \dots\ laudis}, the sacrifice
of praise, ``for the increase of our service,'' and asks God to bring to
completion what he has given \latin{immeritis}, to the undeserving. Like the
Sunday's Collect it was composed for another Mass and is said here because the
Sunday is commemorated. Its closing words are said aloud to begin the Preface.

\subsection{Preface}
\label{sec:element-preface}

The feast has no Preface of its own, and the two Prefaces that the rubrics
admit are set out with the appointed texts. Either of them ends with the
praise of the angels and leads into the \latin{Sanctus}.

\subsection{Communion: three beatitudes again}
\label{sec:element-communion}

The Communion repeats the sixth, seventh and eighth beatitudes, the clean of
heart, the peacemakers and those persecuted for justice, word for word.
Fromage gives a reason for the choice of these three, which the second
interpretation below quotes. Schuster calls the Beatitudes ``as it were the
decalogue of happiness'' and paraphrases the three: the clean of heart shall
see God, the peaceable shall be known as true sons of God, and the persecuted
will receive endless life and joy for what they gave up (p.~211).

\subsection{Postcommunion: veneration and protection}
\label{sec:element-postcommunion}

The Postcommunion asks that the faithful peoples ever rejoice in the
veneration of all the saints and be defended by their unceasing prayer. It
stands in Ménard's Gregorian Sacramentary as a second Postcommunion and in the
Monte Suavicino book as the only one. Schuster dwells on its name for the
Church on earth, \latin{fideles populi}: ``the people journeying towards
eternity with the eyes, and in the light of faith'' (p.~211). St Bernard asks
what our praise does for the saints, and answers that they need none of our
honours: \latin{quod eorum memoriam veneramur, nostra interest, non ipsorum},
``that we venerate their memory concerns us, not them.'' The memory, he says,
kindles in us a desire first for their company, to be gathered and to rejoice
together in the communion of all the saints (sermo~5, n.~5).

\subsection{Commemoration Postcommunion: divine things and human dangers}
\label{sec:element-commemoration-postcommunion}

The Sunday's Postcommunion asks that those to whom God gives to rejoice in
sharing divine things, \latin{divina \dots\ participatione}, he not allow to
be subject to human dangers. It is the last proper prayer of the Mass, said
after the feast's Postcommunion, so that the feast's prayer to be defended by
the saints' prayer is followed by the Sunday's prayer against human dangers.
